\chapter{相关工作与技术基础}
\label{chap:background}

本章对与本文研究密切相关的三大方向进行系统梳理：检索增强生成与对话检索的技术演进（第~\ref{sec:bg_rag}~节），查询重写方法的研究现状（第~\ref{sec:bg_qr}~节），以及后续章节将用到的大模型训练技术基础（第~\ref{sec:bg_train}~节）。

\section{检索增强生成与对话检索}
\label{sec:bg_rag}

\subsection{RAG 技术路线总览}

RAG 的基本形态由 Lewis 等人~\citep{lewis2020rag} 与 Guu 等人~\citep{guu2020realm} 系统提出，其经典流水线为"查询编码→稠密向量检索→生成"三阶段。近年来 Gao 等人的综述~\citep{gao2024retrievalaugmented} 将 RAG 研究梳理为 Naive RAG、Advanced RAG 与 Modular RAG 三代：Naive RAG 采用"一次检索，单次生成"的朴素架构；Advanced RAG 在检索前后分别引入查询变换与重排序等模块，提升召回与精度；Modular RAG 则通过 Agent 化、自适应检索等机制赋予系统动态决策能力，本文所聚焦的查询重写即属于 Advanced RAG 中最关键的检索前置模块。

\subsection{稠密检索器}

稠密向量检索是 RAG 系统的事实标准检索后端。DPR~\citep{karpukhin-etal-2020-dense} 首次把双塔 BERT~\citep{devlin2019bert} 编码器用于开放域 QA 检索，使用对比学习从 QA 对中学习相关性信号；Contriever~\citep{izacard2022contriever} 进一步探索了无监督对比预训练，展示了在大规模无标注语料上预训练检索器的可行性。BGE~\citep{xiao2023cpack} 通过大规模多阶段训练构建了社区中使用最广的开源嵌入模型系列。Dragon+~\citep{lin2023dragon} 通过多源数据增强显著提升了检索器的跨域泛化能力；基于 Dragon+ 进一步微调的 Dragon-ChatQA / Dragon-DocChat~\citep{Liu2024ChatQASG} 则把对话历史直接拼接进查询端，成为端到端对话检索的代表。Jasper \& Stella~\citep{Zhang2024JasperAS} 和 gte-ModernBERT~\citep{warner2024modernbert} 等最新工作在精度与推理效率上进一步推高了稠密检索器的上限。

这些进展使得\textbf{在单轮查询意图明确的条件下}，稠密检索器已能够逼近人工上界；然而在多轮对话检索场景中，查询意图本身是残缺、动态且话题漂移的，检索器直接面对原始轮次时性能急剧下降。这一差距正是查询重写等上游模块存在的根本原因。

\subsection{多轮对话检索的数据资源}

围绕多轮对话检索已有一批重要的公开基准，可分为四大类：
\begin{itemize}
    \item \textbf{单文档对话问答}：CoQA~\citep{reddy-etal-2019-coqa}、QuAC~\citep{choi-etal-2018-quac}，专注指代与省略现象下的阅读理解；
    \item \textbf{开放域对话检索}：QReCC~\citep{anantha-etal-2021-qrecc}、OR-QuAC~\citep{qu2020orquac}、TopiOCQA~\citep{adlakha-etal-2022-topiocqa}，要求模型从大规模语料库中检索证据；QReCC 与 TopiOCQA 额外提供了人工重写的独立查询，是查询重写研究的主要训练/评测资源；TopiOCQA 还显式引入了话题切换标注；
    \item \textbf{文档引导型对话}：Doc2Dial~\citep{feng-etal-2020-doc2dial}，面向任务导向型、结构化文档上的多轮问答；
    \item \textbf{大规模合成基准}：CORAL~\citep{coral2025} 等利用大模型自动合成大规模对话 RAG 评测集。
\end{itemize}

本文选取 QReCC、QuAC、Doc2Dial、TopiOCQA 作为主评测基准，并额外使用实验室合作构建的 MTR-Bench~\citep{mtrsuite2025} 作为困难外部评测集，以检验方法在更极端话题切换与更长回复条件下的鲁棒性。

\section{查询重写方法的研究现状}
\label{sec:bg_qr}

查询重写的目标是将上下文依赖的对话查询 $q_i$ 改写为独立于历史、语义自包含的新查询 $q_i^\star$，使其可直接输入单轮检索器。按技术路线划分，已有工作大致可分为三类。

\subsection{监督式重写}

监督式重写以人工标注的"原问题—独立问题"配对为训练数据，最大化模型输出与参考重写之间的似然。Elgohary 等人~\citep{elgohary2019canyouunpack} 构建的 CANARD 数据集是该方向的基础资源，其上的 Seq2Seq 重写器长期作为基线。QuReTeC~\citep{voskarides2020quretec} 将重写转化为"词项选择"问题，避免了自由文本生成的稳定性问题；其后 Lin 等人~\citep{lin2020conversational}、Vakulenko 等人~\citep{vakulenko2021question} 基于 T5~\citep{raffel2020t5}、GPT-2 开展了大量端到端微调工作。Yu 等人~\citep{yu2020fewshotgenerative} 探索了小样本条件下的生成式重写。

监督式方法的优势在于训练稳定、输出可控，但存在两大根本问题：其一，人工重写标注稀缺昂贵；其二，优化目标是"像参考重写"而非"被检索器召回"，与下游目标存在系统性偏差。本文第~\ref{chap:prw}~章所提出的 PRW 框架在监督式重写的基础上引入了两阶段解耦设计，但不再把参考重写当成唯一训练信号。

\subsection{零样本大语言模型重写}

利用 GPT-3.5/4、Qwen 等大模型直接完成对话重写是近两年的主流路线之一。LLM4CS~\citep{mao2023llm4cs} 系统研究了不同 prompt 与聚合策略对重写效果的影响，提出多路径思维链与答案生成辅助重写的方案；Ye 等人~\citep{ye2023enhancing} 转而强调"信息量更丰富的重写"；Mao 等人~\citep{mao2023search} 提出面向搜索的结构化编辑框架。

大模型重写的明显优势是无需训练、跨领域迁移强；然而其推理成本高、延迟大、对提示词敏感，并且与前面提到的两点问题重叠——不对历史剪枝、目标未对齐检索器。这促使本文在"大模型基底 + 轻量微调"的方向进一步探索。

\subsection{检索反馈驱动的重写对齐}

为消除重写目标与检索偏好之间的失配，一批工作尝试把下游检索结果作为训练信号反向优化重写器。ConvGQR~\citep{mo2023convgqr} 提出生成式重写与潜在答案生成的协同方案；CONQRR~\citep{wu2022conqrr} 首次把检索器的 Recall 作为奖励，用 REINFORCE 对 T5 重写器进行强化学习微调；IterCQR~\citep{jang2024itercqr} 进一步把检索反馈嵌入到迭代自训练循环。Kim 等人~\citep{kim2022saving} 则指出稠密检索器在对话场景下存在对"捷径词汇"的依赖，为重写目标设计提供了参考。Mo 等人~\citep{mo2024convsar} 近期的会话搜索综述系统地回顾了这一分支。

现有工作的共性不足是训练不稳定且算力成本较高。REINFORCE/PPO~\citep{schulman2017ppo} 系方法对奖励方差敏感；迭代自训练易受早期错误放大。直接偏好优化（DPO）~\citep{rafailov2023direct} 把偏好学习问题转化为等价的分类目标，避免显式奖励模型与策略梯度，在稳定性与计算成本上显著优于传统 RLHF~\citep{ouyang2022traininglanguagemodelsfollow}。本文第~\ref{chap:rfr}~章的 RFR 方案即以 DPO 为主框架，把检索召回差异作为偏好信号，形成无需额外人工标注、训练稳定、可在中等算力下完成的新型检索反馈对齐方法。

\section{大模型训练技术基础}
\label{sec:bg_train}

\subsection{有监督微调与参数高效训练}

本文的 PRW 与 RFR 均基于 Qwen2.5~\citep{qwen2025qwen25technicalreport} 系列开源基座模型。有监督微调（Supervised Fine-Tuning, SFT）通过最大化监督数据上条件似然 $\sum_{(x,y)\in\mathcal{D}} \log p_\theta(y\mid x)$ 来让模型适配下游任务。为降低显存开销与训练成本，本文实验统一采用 LoRA~\citep{lora2022} 进行参数高效微调：冻结基座模型权重，在注意力投影层注入低秩可训练适配器矩阵，仅更新约 0.5\%--2\% 的参数即可达到全参数微调的 90\% 以上性能。

\subsection{直接偏好优化（DPO）}

DPO~\citep{rafailov2023direct} 是本文 RFR 方案的核心训练算法。给定偏好对 $(x, y^+, y^-)$，DPO 定义以下分类式目标：
\begin{equation}
\mathcal{L}_{\text{DPO}}(\theta) = - \mathbb{E}_{(x, y^+, y^-)}\Big[\log \sigma\big(\beta \log\tfrac{\pi_\theta(y^+\mid x)}{\pi_{\text{ref}}(y^+\mid x)} - \beta \log\tfrac{\pi_\theta(y^-\mid x)}{\pi_{\text{ref}}(y^-\mid x)}\big)\Big],
\label{eq:dpo_bg}
\end{equation}
其中 $\pi_\theta$ 为待优化策略，$\pi_{\text{ref}}$ 为参考策略（通常取 SFT 后的模型），$\beta$ 控制偏好信号强度。相比 PPO，DPO 不需要显式训练奖励模型，也不需要在线采样，仅需一次离线偏好对构造即可完成优化，训练稳定性与工程复杂度显著优化。

\subsection{策略梯度与 PPO（对照）}

作为本文 RFR 的对照方法之一，PPO~\citep{schulman2017ppo} 通过裁剪后的替代目标函数在保持策略不过度偏离参考策略的条件下最大化期望奖励。其主要挑战在于奖励函数的稀疏性与方差，以及在线采样的计算开销。RLHF~\citep{ouyang2022traininglanguagemodelsfollow} 在 InstructGPT 工作中将 PPO 与奖励模型结合，但在查询重写这类奖励信号稀疏、文本较短的场景下训练极为不稳定，这也是本文最终选用 DPO 而非 PPO 作为主路线的经验依据。

\section{本章小结}

本章从 RAG 与稠密检索、对话检索数据资源、查询重写三条技术路线、以及大模型训练方法基础四个方面系统梳理了与本文相关的研究现状，明确了本文方法的技术出发点。第~\ref{chap:framework}~章将在上述知识基础上，给出本文查询重写任务的形式化定义与三阶段框架设计。
